\scanpage{14}
\illegible{2}\unsure{toute} \inserted{fonction pol. de \unsure{degré} \unsure{\illegible{1}|$\leq2$}} $q$ qui s'annu\unsure{le} sur cinq sommets consécutifs s'annule sur tous les sommets (si $\alpha(\alpha+1)$ inversible) et est en fait un multiple de la forme \deleted{quadratiq} invariante principale.

Sur l'espace des translations, la forme quad. principale est
\[
\deleted{\unsure{q_0(x,y)|q(x,y)}}\quad q_{\deleted{\unsure{0}}\unsure{0|\Phi}}(x,y)=x^2+y^2-\alpha xy
\tag*{(3)}
\]
\deleted{qui est dég.} dont le discriminant est
\[
\delta_0=\alpha^2-4=(\alpha-2)(\alpha+2)\tag*{(4)}
\]
inv. sss $\alpha\ne\unsure{\pm|+}2\deleted{\unsure{0}}$ fibre par fibre i.e. sss on est dans cas I.

Quant à $q_0$ elle-même, s'écrit comme forme en $3$ variables $x,y,z$
\[
q_{\unsure{0|\Phi}}(x,y,z)=x^2+y^2\deleted{\unsure{xy}}+0z^2-yz-\unsure{z}x-\alpha xy
\tag*{(2 bis)}
\]
de matrice adjointe
\[
\begin{pmatrix}2&-\alpha&-1\\-\alpha&2&-1\\-1&-1&0\end{pmatrix}
\tag*{(5)}
\]
de discriminant
\[
\delta=\deleted{\unsure{4a+}}-2\alpha\deleted{\unsure{+}}\unsure{-|+}2-2=-2(\alpha+2)
\tag*{(6)}
\]
de discriminant divisé
\[
\delta'=-(\alpha+2)\tag*{(6$'$)}
\]
qui est nul sss $\alpha=-2$, inv. sss $\alpha\ne-\unsure{2|1}$ sur chaque fibre (i.e. si on n'est pas dans cas $\overline{\underline{\mathrm{III}}}$ sur aucune fibre...)
% archive sheet number: 13
